\section{Konstruksi Hamiltonian dan Inisialisasi Kuantum}
Hamiltonian total dalam penelitian ini dikonstruksi secara sistematis melalui penggabungan antara fungsi biaya portofolio ($H_{cost}$) dan fungsi penalti kendala ($H_{penalty}$) guna memandu proses pencarian solusi yang valid secara ekonomi dan teknis. Suku $H_{cost}$ merepresentasikan lanskap risiko dan ekspektasi imbal hasil yang telah dipetakan secara akurat dari parameter interaksi $J_{ij}$ serta bias individu $h_i$ yang diperoleh pada tahapan analisis ekonofisika sebelumnya. Sementara itu, suku $H_{penalty}$ berfungsi sebagai mekanisme kontrol yang ketat untuk menegakkan batasan \textit{Hamming Weight}, yang memaksa sistem kuantum agar hanya memilih tepat sejumlah $K$ aset ke dalam konfigurasi portofolio final. Penggunaan konstanta penalti ($\lambda$) yang memiliki nilai cukup besar menjamin bahwa setiap status basis qubit yang tidak memenuhi kendala jumlah aset akan memiliki tingkat energi yang sangat tinggi, sehingga secara alami akan dieliminasi oleh algoritma selama proses minimisasi energi berlangsung.

Guna mengeksplorasi ruang solusi yang sangat luas, penelitian ini mengadopsi arsitektur \textit{Hardware-Efficient Ansatz} (HEA) yang dirancang secara khusus untuk meminimalkan kedalaman sirkuit serta dampak negatif dari efek derau pada perangkat kuantum era NISQ. Ansatz ini secara struktural terdiri dari lapisan gerbang rotasi qubit yang dapat diparameterisasi ($\theta$) secara bebas, serta lapisan gerbang keterbelitan (\textit{entanglement}) yang berfungsi untuk memodelkan struktur korelasi antar aset finansial. Inisialisasi dari setiap parameter $\theta$ dilakukan secara cerdas dengan memanfaatkan hasil ekuilibrium Nash yang diperoleh dari kerangka kerja \textit{Exact Potential Game} sebagai titik awal pencarian atau \textit{warm start}. Setiap aset tunggal yang terdaftar dalam indeks saham LQ45 dipetakan secara satu-ke-satu ke dalam sebuah qubit, di mana status fisik qubit $|1\rangle$ memberikan interpretasi bahwa aset tersebut dipilih untuk dimasukkan ke dalam keranjang portofolio. Operator biner yang terdapat pada formulasi Hamiltonian Ising kemudian dipromosikan menjadi operator \textit{Pauli-Z} ($\hat{Z}$) yang dapat dievaluasi dan diolah secara langsung oleh sirkuit kuantum variasional melalui operasi gerbang logika.
